\documentclass[10pt,a4paper]{amsart}
\usepackage[margin=19mm]{geometry}
\usepackage{amsmath,amssymb,mathtools}
\usepackage[scr=rsfs,cal=euler]{mathalfa}
\usepackage{xfrac}
\usepackage[unicode,hidelinks,bookmarks=false]{hyperref}
\newcommand{\ie}{i.e.\ }
\newcommand{\cf}{cf.\ }
\newcommand{\resp}{respectively\ }
\newcommand{\Subsection}{\subsection}
\providecommand{\nobreakdash}{\nobreak\hbox{-}}
\setlength{\emergencystretch}{2em}
\multlinegap0pt
\usepackage{bm}
\usepackage{bbm}
\usepackage{dsfont}
\usepackage{xspace}
\usepackage{cancel}
\usepackage{mathtools}
\usepackage{amsthm}
\makeatletter
\@ifundefined{theorem}{\newtheorem{theorem}{Theorem}}{}
\@ifundefined{conjecture}{\newtheorem{conjecture}[theorem]{Conjecture}}{}
\@ifundefined{definition}{\newtheorem{definition}[theorem]{Definition}}{}
\@ifundefined{proposition}{\newtheorem{proposition}[theorem]{Proposition}}{}
\@ifundefined{remark}{\newtheorem{remark}[theorem]{Remark}}{}
\makeatother
\title[Equivalence of linear and trilinear Kakeya conjectures in th]{Equivalence of linear and trilinear Kakeya conjectures in three dimensions}
\author{Cristian Rios, Eric T. Sawyer}
\date{Source: arXiv:2506.21315v4 (CC BY 4.0)}
\begin{document}
\maketitle

\noindent\emph{Source and licence.} Equivalence of linear and trilinear Kakeya conjectures in three dimensions. Version arXiv:2506.21315v4 (CC BY 4.0), distributed under the Creative Commons Attribution 4.0 International licence, as recorded in the arXiv licence metadata for this exact version and shown on the abstract page https://arxiv.org/abs/2506.21315v4. Full rights notice: licenses/arxiv-2506.21315-CC-BY-4.0.txt.\par\smallskip

\noindent\textit{Editorial scope.} Counted unit: the Proposition whose source label is \texttt{Kak tri equiv} in the article (source0.tex lines 864--869, source label \texttt{Kak tri equiv}), delivered together with the article's own complete original proof of it (source0.tex lines 871--1046, one \texttt{proof} environment of 176 lines). No mathematics is written, repaired or re-derived: statement and proof are the article's own text line for line. Required context retained uncounted: Kak dual (lines 114--126); linear Kakeya (lines 135--140); tri Kak dual (lines 159--174); trilinear Kakeya (lines 176--181); ssFsfec (lines 418--423); Kak Four (lines 449--453); square point (lines 510--515); interp (lines 561--577); single tri Four (lines 595--615); ssFsftec (lines 617--623); Kak equiv (lines 772--776); rap (lines 782--799); namely (lines 2296--2300); param (lines 2320--2331). Every definition, display and label of those lines is retained unchanged. Omitted from this package and not redistributed: 1--113, 127--134, 141--158, 175--175, 182--417, 424--448, 454--509, 516--560, 578--594, 616--616, 624--771, 777--781, 800--863, 870--870, 1047--2295, 2301--2319, 2332--2473. Nothing inside a retained excerpt is omitted or altered: every retained body is delivered exactly as the article's lines are written, so no omission receipt of any kind accompanies this package. Citations: the retained excerpts carry no citation command at all, so no bibliography entry of the article is transcribed and no reference is redistributed. Reference adaptation: none was needed and none was made; every reference inside the delivered unit resolves inside it, so no unresolved reference or citation is produced. Printed slips kept literal: no printed word, symbol, hypothesis or number of the counted statement or of its proof was altered. Layout and class: the article's own class is \texttt{amsart} with packages including amssymb, amsfonts, geometry, graphicx, amsmath; the wrapper is the lane's pinned \texttt{amsart} editorial wrapper at 10pt on A4, which restates the theorem-like environments the retained text needs and adds the article's own macro definitions that the retained text needs (none), with extra wrapper packages (bm, bbm, dsfont, xspace, cancel, mathtools). The delivered numbering is the unit's own. Assets: no retained span contains a figure environment, a graphics-inclusion command, a PDF-page inclusion, a table or a TikZ picture; no image file is copied into this package, the package declares no asset dependency and no image file is redistributed with it. Licence: version arXiv:2506.21315v4 of this article is distributed under the Creative Commons Attribution 4.0 International licence, as recorded in the arXiv licence metadata for this exact version and shown on the abstract page https://arxiv.org/abs/2506.21315v4. A full rights notice is delivered at licenses/arxiv-2506.21315-CC-BY-4.0.txt. Characters: every retained excerpt is pure ASCII, so the delivered engine, XeLaTeX with its default Latin Modern text font, reports no missing character.
\begin{definition}
Let $0<\varepsilon<1$. We say the statement $\mathcal{K}^{\ast}\left(
\otimes_{1}L^{\infty}\rightarrow L^{\frac{3}{2}};\varepsilon\right)  $ holds
if there is a positive constant $C_{\varepsilon}$ such that%
\begin{align}
&  \left\Vert \sum_{T\in\mathbb{T}}\mathbf{1}_{T}\right\Vert _{L^{\frac{3}{2}%
}\left(  \mathbb{R}^{3}\right)  }\leq C_{\varepsilon}\delta^{-\varepsilon
},\label{Kak dual}\\
&  \text{for all families }\mathbb{T}\text{ of }\delta\text{-separated }%
\delta\text{-tubes in }\mathbb{R}^{3}\text{ and }0<\delta<1.\nonumber
\end{align}

\end{definition}

\begin{conjecture}
[Dual form of the Kakeya maximal operator conjecture in $\mathbb{R}^{3}$%
]\label{linear Kakeya}The statement $\mathcal{K}^{\ast}\left(  \otimes
_{1}L^{\infty}\rightarrow L^{\frac{3}{2}};\varepsilon\right)  $ holds for all
$0<\varepsilon<1$.
\end{conjecture}

\begin{definition}
Let $0<\varepsilon,\nu<1$. We say the statement $\mathcal{K}%
_{\operatorname*{disj}\nu}^{\ast}\left(  \otimes_{3}L^{\infty}\rightarrow
L^{\frac{1}{2}};\varepsilon\right)  $ holds if there is a positive constant
$C_{\varepsilon,\nu}$ such that%
\begin{align}
&  \left\Vert \prod_{k=1}^{3}\left(  \sum_{T_{k}\in\mathbb{T}_{k}}%
\mathbf{1}_{T_{k}}\right)  \right\Vert _{L^{\frac{1}{2}}\left(  \mathbb{R}%
^{3}\right)  }\leq C_{\varepsilon,\nu}\delta^{-\varepsilon}%
,\label{tri Kak dual}\\
&  \text{for all }\nu\text{\emph{-disjoint }families }\mathbb{T}_{k}\text{ of
}\delta\text{-separated }\delta\text{-tubes in }\mathbb{R}^{3}\text{ and
}0<\delta\leq\nu.\nonumber
\end{align}

\end{definition}

\begin{conjecture}
[The disjoint trilinear dual form of the Kakeya maximal operator conjecture in
$\mathbb{R}^{3}$]\label{trilinear Kakeya}The statement $\mathcal{K}%
_{\operatorname*{disj}\nu}^{\ast}\left(  \otimes_{3}L^{\infty}\rightarrow
L^{\frac{1}{2}};\varepsilon\right)  $ holds for all $0<\varepsilon,\nu<1$.
\end{conjecture}

\begin{conjecture}
[modulated single scale Fourier square function extension conjecture]%
\label{ssFsfec}The statement $\mathcal{E}^{\operatorname*{square}}\left(
\otimes_{1}L^{\infty}\rightarrow L^{q};\varepsilon\right)  $ holds for all
$0<\varepsilon<1$ and $q>3$.
\end{conjecture}

\begin{proposition}
\label{Kak Four}The modulated single scale Fourier square function extension
Conjecture \ref{ssFsfec} implies the dual form of the Kakeya maximal operator
Conjecture \ref{linear Kakeya}.
\end{proposition}

\begin{equation}
\mathcal{S}\mathbb{T}\left(  2^{s}\times2^{s}\times2^{2s}\right)  \left(
\xi\right)  \leq C2^{2s}\mathcal{S}_{\operatorname*{Fourier}}^{s,\mathbf{u}%
}f\left(  \xi\right)  ,\ \ \ \ \ \text{for all }\xi\in\mathbb{R}^{3},
\label{square point}%
\end{equation}

\begin{align}
&  \int_{\mathbb{R}^{3}}\left(  \sum_{T\in\mathbb{T}\left(  2^{-s}\times
2^{-s}\times1\right)  }\mathbf{1}_{T}\left(  \xi\right)  \right)  ^{\frac
{3}{2}}d\xi=\left\{  \int_{\left\{  \sum_{T\in\mathbb{T}}\mathbf{1}_{T}%
\leq1\right\}  }+\int_{\left\{  \sum_{T\in\mathbb{T}}\mathbf{1}_{T}>1\right\}
}\right\}  \left(  \sum_{T\in\mathbb{T}\left(  2^{-s}\times2^{-s}%
\times1\right)  }\mathbf{1}_{T}\left(  \xi\right)  \right)  ^{\frac{3}{2}}%
d\xi\label{interp}\\
&  \leq\int_{\left\{  \sum_{T\in\mathbb{T}}\mathbf{1}_{T}\leq1\right\}
}\left(  \sum_{T\in\mathbb{T}\left(  2^{-s}\times2^{-s}\times1\right)
}\mathbf{1}_{T}\left(  \xi\right)  \right)  d\xi+\int_{\left\{  \sum
_{T\in\mathbb{T}}\mathbf{1}_{T}>1\right\}  }\left(  \sum_{T\in\mathbb{T}%
\left(  2^{-s}\times2^{-s}\times1\right)  }\mathbf{1}_{T}\left(  \xi\right)
\right)  ^{p_{\varepsilon}}d\xi\nonumber\\
&  \lesssim1+C_{p_{\varepsilon}}2^{\left(  4p_{\varepsilon}-6\right)
s}=1+C_{p_{\varepsilon}}2^{4\varepsilon s}.\nonumber
\end{align}

\begin{definition}
Let $1<q<\infty$ and $0<\varepsilon,\nu<1$. We say the statement
$\mathcal{A}_{\operatorname*{disj}\nu}^{\operatorname*{square}}\left(
\otimes_{3}L^{\infty}\rightarrow L^{\frac{q}{3}};\varepsilon\right)  $ holds
if there is a positive constant $C_{q,\varepsilon,\nu}$ depending only on $q$,
$\varepsilon$ and $\nu$, such that,%
\begin{align}
&  \left\Vert \mathcal{S}_{\operatorname*{Fourier}}^{s,\mathbf{u}_{1}}%
f_{1}\ \mathcal{S}_{\operatorname*{Fourier}}^{s,\mathbf{u}_{2}}f_{2}%
\ \mathcal{S}_{\operatorname*{Fourier}}^{s,\mathbf{u}_{3}}f_{3}\right\Vert
_{L^{\frac{q}{3}}\left(  \mathbb{R}^{3}\right)  }\leq C_{q,\varepsilon,\nu
}2^{\varepsilon s}\left\Vert f_{1}\right\Vert _{L^{\infty}\left(  U\right)
}\left\Vert f_{2}\right\Vert _{L^{\infty}\left(  U\right)  }\left\Vert
f_{3}\right\Vert _{L^{\infty}\left(  U\right)  }\ ,\label{single tri Four}\\
\text{for all }s  &  \in\mathbb{N}\text{ with }2^{-s}\leq\nu\text{, all }%
f_{k}\in L^{\infty}\left(  U_{k}\right)  \text{, all sequences }\mathbf{u}%
_{k}\in\mathcal{V}\text{, and all }\nu\text{-disjoint triples }\left(
U_{1},U_{2},U_{3}\right)  \subset U^{3}.\nonumber
\end{align}

\end{definition}

\begin{conjecture}
[modulated single scale Fourier square function disjoint trilinear extension
conjecture]\label{ssFsftec}For every $q>3$ there is $0<\nu<1$ such that the
statement $\mathcal{A}_{\operatorname*{disj}\nu}^{\operatorname*{square}%
}\left(  \otimes_{3}L^{\infty}\rightarrow L^{\frac{q}{3}};\varepsilon\right)
$ holds for all $0<\varepsilon<1$.
\end{conjecture}

\begin{proposition}
\label{Kak equiv}The dual form of the Kakeya maximal operator Conjecture
\ref{linear Kakeya} holds \emph{if and only if} the single scale modulated
Fourier square function extension Conjecture \ref{ssFsfec} holds.
\end{proposition}

\begin{align}
&  \left\Vert \mathcal{S}_{\operatorname*{Fourier}}^{s,\mathbf{u}}f\right\Vert
_{L^{3}\left(  \mathbb{R}^{3}\right)  }=\left\Vert \left(  \sum_{I\in
\mathcal{G}_{s}\left[  U\right]  }\left\vert \tau_{u_{I}}\widehat{\Phi_{\ast
}\bigtriangleup_{I}f}\left(  \xi\right)  \right\vert ^{2}\right)  ^{\frac
{1}{2}}\right\Vert _{L^{3}\left(  \mathbb{R}^{3}\right)  }\label{rap}\\
&  \lesssim2^{-2s}\left(  \int_{\mathbb{R}^{3}}\left(  \sum_{I\in
\mathcal{G}_{s}\left[  U\right]  }\left\vert \mathbf{1}_{\tau_{u_{I}}%
\widehat{I}}\left(  \xi\right)  \right\vert ^{2}\right)  ^{\frac{3}{2}}%
d\xi\right)  ^{\frac{1}{3}}\left\Vert f\right\Vert _{L^{\infty}\left(
U\right)  }\nonumber\\
&  \ \ \ \ \ \ \ \ \ \ \ \ \ \ \ \ \ \ \ \ \ \ \ \ \ \ \ \ \ \ +\sum
_{m=1}^{\infty}a_{m}2^{-2s}\left(  \int_{\mathbb{R}^{3}}\left(  \sum
_{I\in\mathcal{G}_{s}\left[  U\right]  }\left\vert \mathbf{1}_{2^{m}%
\tau_{u_{I}}\widehat{I}}\left(  \xi\right)  \right\vert ^{2}\right)
^{\frac{3}{2}}d\xi\right)  ^{\frac{1}{3}}\left\Vert f\right\Vert _{L^{\infty
}\left(  U\right)  }\nonumber
\end{align}

\begin{equation}
2^{-\left(  q-3\right)  \frac{1}{q}\lambda}+2^{-\left(  q-\frac{3}{2}\right)
\frac{2}{q}\lambda}+2^{-\left(  2-\frac{4}{q}\right)  \lambda}<\frac{1}{2}.
\label{namely}%
\end{equation}

\begin{remark}
\label{param}The inequality (\ref{namely}) holds if
\[
\left(  1-\frac{3}{q}\right)  \lambda>3\text{ and }\left(  1-\frac{2}%
{q}\right)  \lambda>1\text{, in particular if }\lambda>\frac{3q}{q-3}.
\]
Thus altogether we need to take%
\[
\lambda>\max\left\{  \frac{3q}{q-3},\log_{2}\frac{2^{10}}{\nu_{q}}\right\}  .
\]

\end{remark}

\begin{proposition}
\label{Kak tri equiv}The disjoint trilinear dual form of the Kakeya maximal
operator Conjecture \ref{trilinear Kakeya} holds \emph{if and only if} the
modulated single scale disjoint\emph{\ }trilinear Fourier square function
Conjecture \ref{ssFsftec} holds.
\end{proposition}

\begin{proof}
The `if' assertion is similar to the argument proving (\ref{Kak dual}) in
Proposition \ref{Kak Four}, but with trilinearity in place of linearity.
Indeed, given $0<\varepsilon,\nu<1$, we must show that $\mathcal{K}%
_{\operatorname*{disj}\nu}^{\ast}\left(  \otimes_{3}L^{\infty}\rightarrow
L^{\frac{3}{2}};\varepsilon\right)  $ holds. Since $\mathcal{K}%
_{\operatorname*{disj}\nu^{\prime}}^{\ast}\left(  \otimes_{3}L^{\infty
}\rightarrow L^{\frac{3}{2}};\varepsilon^{\prime}\right)  $ implies
$\mathcal{K}_{\operatorname*{disj}\nu}^{\ast}\left(  \otimes_{3}L^{\infty
}\rightarrow L^{\frac{3}{2}};\varepsilon\right)  $ for $\varepsilon^{\prime
}\leq\varepsilon$ and $\nu^{\prime}\leq\nu$, it is enough to show that there
exists $\nu^{\prime}\in\left(  0,\nu\right)  $ such that $\mathcal{K}%
_{\operatorname*{disj}\nu^{\prime}}^{\ast}\left(  \otimes_{3}L^{\infty
}\rightarrow L^{\frac{3}{2}};\varepsilon^{\prime}\right)  $ holds for all
$0<\varepsilon^{\prime}<1 $. For this we fix $q_{\varepsilon}\in\left(
3,3+\varepsilon\right)  $. Then by Conjecture \ref{ssFsfec} there is
$\nu^{\prime}>0$ depending on $q_{\varepsilon}$ such that
(\ref{single tri Four}) holds for all $0<\varepsilon^{\prime}<1$, and in light
of Remark \ref{param}, we may assume $\nu^{\prime}<\nu$. Thus if $\left(
\mathbb{T}_{1},\mathbb{T}_{2},\mathbb{T}_{3}\right)  $ is a $\nu^{\prime}%
$-disjoint triple of families of $2^{-s}$-separated $2^{-s}$-tubes, and if
$\left(  U_{1},U_{2},U_{3}\right)  $ is a $\nu^{\prime}$-disjoint triple of
squares in $\mathcal{G}\left[  U\right]  $ such that the orientations of the
tubes in $\mathbb{T}_{k}$ are normal to some point on the surface $\Phi\left(
U_{k}\right)  $, we have%
\begin{align*}
&  \left(  \int_{\mathbb{R}^{3}}\left(  \mathcal{S}_{\operatorname*{Fourier}%
}^{s,\mathbf{u}_{1}}f_{1}\left(  \xi\right)  \ \mathcal{S}%
_{\operatorname*{Fourier}}^{s,\mathbf{u}_{2}}f_{2}\left(  \xi\right)
\ \mathcal{S}_{\operatorname*{Fourier}}^{s,\mathbf{u}_{3}}f_{3}\left(
\xi\right)  \right)  ^{\frac{q_{\varepsilon}}{3}}\ d\xi\right)  ^{\frac
{3}{q_{\varepsilon}}}\\
&  \ \ \ \ \ \ \ \ \ \ \ \ \ \ \ \ \ \ \ \ \lesssim2^{\varepsilon^{\prime}%
s}\left\Vert f_{1}\right\Vert _{L^{\infty}\left(  U_{1}\right)  }\left\Vert
f_{2}\right\Vert _{L^{\infty}\left(  U_{2}\right)  }\left\Vert f_{3}%
\right\Vert _{L^{\infty}\left(  U_{3}\right)  }\ ,
\end{align*}
with $f_{k}\equiv\sum_{I\in G_{s}\left[  U_{k}\right]  }\bigtriangleup_{I}1$.
Then using the pointwise square function estimate (\ref{square point}),
\[
\mathcal{S}\mathbb{T}\left(  2^{s}\times2^{s}\times2^{2s}\right)  \left(
\xi\right)  \leq C2^{2s}\mathcal{S}_{\operatorname*{Fourier}}^{s,\mathbf{u}%
}f\left(  \xi\right)  ,
\]
we obtain from this that,%
\begin{align*}
&  2^{-6s}\left\Vert \prod_{k=1}^{3}\mathcal{S}\mathbb{T}_{k}\left(
2^{s}\times2^{s}\times2^{2s}\right)  \right\Vert _{L^{\frac{q_{\varepsilon}%
}{3}}\left(  \mathbb{R}^{3}\right)  }=\left\Vert \prod_{k=1}^{3}%
2^{-2s}\mathcal{S}\mathbb{T}_{k}\left(  2^{s}\times2^{s}\times2^{2s}\right)
\right\Vert _{L^{\frac{q_{\varepsilon}}{3}}\left(  \mathbb{R}^{3}\right)  }\\
&  \lesssim\left\Vert \prod_{k=1}^{3}\mathcal{S}_{\operatorname*{Fourier}%
}^{s,\mathbf{u}_{k}}f_{k}\left(  \xi\right)  \right\Vert _{L^{\frac{q_{3}}{3}%
}\left(  \mathbb{R}^{3}\right)  }\lesssim2^{\varepsilon^{\prime}s}\left\Vert
f_{1}\right\Vert _{L^{\infty}}\left\Vert f_{2}\right\Vert _{L^{\infty}%
}\left\Vert f_{3}\right\Vert _{L^{\infty}}\lesssim2^{\varepsilon^{\prime}s}\ ,
\end{align*}
for all $\nu^{\prime}$-disjoint families of $2^{-s}$-separatede $2^{-s}%
$-tubes, and hence by rescaling that%
\begin{align*}
&  \left\Vert \prod_{k=1}^{3}\mathcal{S}\mathbb{T}_{k}\left(  2^{-s}%
\times2^{-s}\times1\right)  \right\Vert _{L^{\frac{q_{\varepsilon}}{3}}\left(
\mathbb{R}^{3}\right)  }=\left(  \int_{\mathbb{R}^{3}}\left(  \prod_{k=1}%
^{3}\mathcal{S}\mathbb{T}_{k}\left(  2^{-s}\times2^{-s}\times1\right)  \left(
\xi\right)  \right)  ^{\frac{q_{\varepsilon}}{3}}d\xi\right)  ^{\frac
{3}{q_{\varepsilon}}}\\
&  =\left(  \int_{\mathbb{R}^{3}}\left(  \prod_{k=1}^{3}\mathcal{S}%
\mathbb{T}_{k}\left(  2^{-s}\times2^{-s}\times1\right)  \left(  \frac
{\xi^{\prime}}{2^{2s}}\right)  \right)  ^{\frac{q_{\varepsilon}}{3}}d\frac
{\xi^{\prime}}{2^{2s}}\right)  ^{\frac{3}{q_{\varepsilon}}}%
\ \ \ \ \ \text{(with }\xi=\frac{\xi^{\prime}}{2^{2s}}\text{)}\\
&  =\left(  \int_{\mathbb{R}^{3}}\left(  \prod_{k=1}^{3}\mathcal{S}%
\mathbb{T}_{k}\left(  2^{s}\times2^{s}\times2^{2s}\right)  \left(  \xi
^{\prime}\right)  \right)  ^{\frac{q_{\varepsilon}}{3}}2^{-6s}d\xi^{\prime
}\right)  ^{\frac{3}{q_{\varepsilon}}}\\
&  =\left(  2^{-6s}\right)  ^{\frac{3}{q_{\varepsilon}}}\left\Vert \prod
_{k=1}^{3}\mathcal{S}\mathbb{T}_{k}\left(  2^{s}\times2^{s}\times
2^{2s}\right)  \right\Vert _{L^{\frac{q_{\varepsilon}}{3}}\left(
\mathbb{R}^{3}\right)  }\lesssim2^{\left(  6-\frac{18}{q_{\varepsilon}%
}\right)  s}2^{\varepsilon^{\prime}s}<2^{\left(  6-\frac{18}{3+\varepsilon
}+\varepsilon^{\prime}\right)  s}=2^{\left(  \frac{6\varepsilon}%
{3+\varepsilon}+\varepsilon^{\prime}\right)  s}<2^{3\varepsilon s},
\end{align*}
if we take $\varepsilon^{\prime}=\varepsilon$, since $q_{\varepsilon
}<3+\varepsilon$.

Arguing as in (\ref{interp}) we now obtain (\ref{tri Kak dual}) for this
choice of $\nu^{\prime}$. Indeed, with
\[
F\left(  \xi\right)  \equiv\prod_{k=1}^{3}\left(  \sum_{T_{k}\in\mathbb{T}%
_{k}\left(  2^{-s}\times2^{-s}\times1\right)  }\mathbf{1}_{T_{k}}\left(
\xi\right)  \right)  ^{\frac{1}{3}},
\]
we have using $q_{\varepsilon}>3$ that
\begin{align*}
&  \int_{\mathbb{R}^{3}}\prod_{k=1}^{3}\left(  \sum_{T_{k}\in\mathbb{T}%
_{k}\left(  2^{-s}\times2^{-s}\times1\right)  }\mathbf{1}_{T_{k}}\left(
\xi\right)  \right)  ^{\frac{1}{2}}d\xi=\int_{\mathbb{R}^{3}}F\left(
\xi\right)  ^{\frac{3}{2}}d\xi\\
&  =\left\{  \int_{\left\{  F\leq1\right\}  }+\int_{\left\{  F>1\right\}
}\right\}  F\left(  \xi\right)  ^{\frac{3}{2}}d\xi\leq\int_{\left\{
F\leq1\right\}  }F\left(  \xi\right)  d\xi+\int_{\left\{  F>1\right\}
}F\left(  \xi\right)  ^{\frac{q_{\varepsilon}}{2}}d\xi\\
&  \lesssim\int_{\left\{  \sum_{T\in\mathbb{T}}\mathbf{1}_{T}\leq1\right\}
}\prod_{k=1}^{3}\left(  \sum_{T_{k}\in\mathbb{T}_{k}\left(  2^{-s}\times
2^{-s}\times1\right)  }\mathbf{1}_{T_{k}}\left(  \xi\right)  \right)
^{\frac{1}{3}}d\xi+\int_{\left\{  \sum_{T\in\mathbb{T}}\mathbf{1}%
_{T}>1\right\}  }\prod_{k=1}^{3}\left(  \sum_{T_{k}\in\mathbb{T}_{k}\left(
2^{-s}\times2^{-s}\times1\right)  }\mathbf{1}_{T_{k}}\left(  \xi\right)
\right)  ^{\frac{q_{\varepsilon}}{6}}d\xi\\
&  \lesssim1+C_{q_{\varepsilon}}^{\frac{q_{\varepsilon}}{3}}2^{\frac
{q_{\varepsilon}}{3}3\varepsilon s}\lesssim1+2^{4\varepsilon s},
\end{align*}
since first,%
\begin{align*}
&  \int_{\left\{  \sum_{T\in\mathbb{T}}\mathbf{1}_{T}\leq1\right\}  }%
\prod_{k=1}^{3}\left(  \sum_{T_{k}\in\mathbb{T}_{k}\left(  2^{-s}\times
2^{-s}\times1\right)  }\mathbf{1}_{T_{k}}\left(  \xi\right)  \right)
^{\frac{1}{3}}d\xi\\
&  \lesssim\prod_{k=1}^{3}\left(  \int_{\mathbb{R}^{3}}\sum_{T_{k}%
\in\mathbb{T}_{k}\left(  2^{-s}\times2^{-s}\times1\right)  }\mathbf{1}_{T_{k}%
}\left(  \xi\right)  d\xi\right)  ^{\frac{1}{3}}\lesssim\prod_{k=1}^{3}\left(
2^{-2s}\#\mathbb{T}_{k}\left(  2^{-s}\times2^{-s}\times1\right)  \right)
^{\frac{1}{3}}\lesssim1,
\end{align*}
because the tubes in $\mathbb{T}_{k}$ are $2^{-s}$-separated, and second since
for $0<\varepsilon<1$, we have $q_{\varepsilon}<3+\varepsilon\leq4$ and,
\[
\int_{\left\{  \sum_{T\in\mathbb{T}}\mathbf{1}_{T}>1\right\}  }\prod_{k=1}%
^{3}\left(  \sum_{T_{k}\in\mathbb{T}_{k}\left(  2^{-s}\times2^{-s}%
\times1\right)  }\mathbf{1}_{T_{k}}\left(  \xi\right)  \right)  ^{\frac
{q_{\varepsilon}}{6}}d\xi\leq\int_{\mathbb{R}^{3}}\prod_{k=1}^{3}%
\mathcal{S}\mathbb{T}_{k}\left(  \xi\right)  ^{\frac{q_{\varepsilon}}{3}}%
d\xi\lesssim\left(  C_{q_{\varepsilon}}2^{3\varepsilon s}\right)
^{\frac{q_{\varepsilon}}{3}}\lesssim2^{4\varepsilon s}.
\]


Conversely, for the `only if' assertion, we use the rapid decay of the Fourier
transform away from the associated tubes, just as in the proof (\ref{rap}) of
the converse assertion of Proposition \ref{Kak equiv}, to obtain the
inequality,%
\begin{align*}
&  \left\Vert \prod_{k=1}^{3}\mathcal{S}_{\operatorname*{Fourier}%
}^{s,\mathbf{u}_{k}}f\right\Vert _{L^{1}\left(  \mathbb{R}^{3}\right)
}\lesssim\left\Vert \left(  \prod_{k=1}^{3}\sum_{m_{k}=0}^{\infty}a_{m_{k}%
}\sum_{I_{k}\in\mathcal{G}_{s}\left[  U_{k}\right]  }\mathbf{1}_{2^{m_{k}}%
\tau_{\mathbf{u}_{k}}\widehat{I}}\right)  ^{\frac{1}{2}}\right\Vert
_{L^{1}\left(  \mathbb{R}^{3}\right)  }\\
&  \lesssim\left\Vert \left(  \sum_{m_{1}=0}^{\infty}\sum_{m_{2}=0}^{\infty
}\sum_{m_{3}=0}^{\infty}a_{m_{1}}a_{m_{2}}a_{m_{3}}\prod_{k=1}^{3}\sum
_{I_{k}\in\mathcal{G}_{s}\left[  U_{k}\right]  }\mathbf{1}_{2^{m_{k}}%
\tau_{\mathbf{u}_{k}}\widehat{I}}\right)  ^{\frac{1}{2}}\right\Vert
_{L^{1}\left(  \mathbb{R}^{3}\right)  }\\
&  \lesssim\left\Vert \sum_{m_{1}=0}^{\infty}\sum_{m_{2}=0}^{\infty}%
\sum_{m_{3}=0}^{\infty}\sqrt{a_{m_{1}}a_{m_{2}}a_{m_{3}}}\prod_{k=1}%
^{3}\left(  \sum_{I_{k}\in\mathcal{G}_{s}\left[  U_{k}\right]  }%
\mathbf{1}_{2^{m_{k}}\tau_{\mathbf{u}_{k}}\widehat{I}}\right)  ^{\frac{1}{2}%
}\right\Vert _{L^{1}\left(  \mathbb{R}^{3}\right)  }\\
&  \lesssim\sum_{m_{1}=0}^{\infty}\sum_{m_{2}=0}^{\infty}\sum_{m_{3}%
=0}^{\infty}\sqrt{a_{m_{1}}a_{m_{2}}a_{m_{3}}}\left\Vert \prod_{k=1}%
^{3}\left(  \sum_{I_{k}\in\mathcal{G}_{s}\left[  U\right]  }\mathbf{1}%
_{2^{m_{k}}\tau_{\mathbf{u}_{k}}\widehat{I}}\right)  ^{\frac{1}{2}}\right\Vert
_{L^{1}\left(  \mathbb{R}^{3}\right)  }\\
&  \lesssim\sum_{m_{1}=0}^{\infty}\sum_{m_{2}=0}^{\infty}\sum_{m_{3}%
=0}^{\infty}\sqrt{a_{m_{1}}a_{m_{2}}a_{m_{3}}}C_{0}^{m_{1}+m_{2}+m_{3}%
}\left\Vert \prod_{k=1}^{3}\left(  \sum_{I_{k}\in\mathcal{G}_{s}\left[
U\right]  }\mathbf{1}_{\tau_{\mathbf{u}_{k}}\widehat{I}}\right)  ^{\frac{1}%
{2}}\right\Vert _{L^{1}\left(  \mathbb{R}^{3}\right)  }\lesssim2^{\varepsilon
s},
\end{align*}
since the coefficients $a_{m}$ are rapidly decreasing, and where again, the
constant $C_{0}$ is related to the geometry of expanded tubes. This shows that
(\ref{single tri Four}) holds for all $\varepsilon,\nu>0$, and completes the
proof of Proposition \ref{Kak tri equiv}.
\end{proof}

\end{document}
